% Overview figure (Figure 1). Panel (a) and the boxed statements are drawn
% here; panels (b) and (c) and the palette are generated by make_figure1.py,
% which also compiles this file. Fonts match the paper: Times text and
% Computer Modern math. Coordinates are in inches.
\documentclass[10pt,border=0pt]{standalone}
\usepackage{times}
\usepackage{amsmath,amssymb}
\usepackage{mathtools}
\usepackage{graphicx}
\usepackage{tikz}
\usetikzlibrary{arrows.meta,calc,shapes.geometric}

\hyphenpenalty=10000 \exhyphenpenalty=10000

\input{palette.tex}
\definecolor{panelbg}{HTML}{F7F8FA}   % panel fill, very pale cool grey
\definecolor{border}{HTML}{C7CDD3}    % every drawn box edge
\definecolor{mathbg}{HTML}{FFFFFF}    % equation boxes

% Three type sizes only, floor 7.5pt against the paper's 10pt body.  The two
% definitions share the header's 9pt rather than taking a fourth size: they are
% what the rest of the figure is written in, so they belong at the top of the
% scale, and the audit stays at seven distinct sizes in the PDF.
\newcommand{\hd}{\fontsize{9}{10.8}\selectfont}    % panel header (a claim)
\newcommand{\eqn}{\fontsize{9}{10.8}\selectfont}   % the two definitions
\newcommand{\bd}{\fontsize{8}{9.6}\selectfont}     % boxed bound
\newcommand{\sm}{\fontsize{7.5}{9}\selectfont}     % chips and secondary notes

% Apply solid red-green-red-green PDF color bands to the term, brace and
% label. The form preserves the original math box and searchable vector text.
% Read endpoints from the same named red/green palette used by the figure.
\newbox\gradientbox
\newbox\gradientformbox
\makeatletter
\def\gradientrgb#1,#2,#3\relax{#1 #2 #3}
\newcommand\gradientmath[1]{\begingroup
  \setbox\gradientbox=\hbox{$\m@th#1$}%
  \edef\gradientwidth{\strip@pt\dimexpr\wd\gradientbox*7200/7227\relax}%
  \convertcolorspec{named}{under}{rgb}\gradientstart
  \convertcolorspec{named}{over}{rgb}\gradientend
  \edef\gradientstart{\expandafter\gradientrgb\gradientstart\relax}%
  \edef\gradientend{\expandafter\gradientrgb\gradientend\relax}%
  \immediate\pdfobj{<< /Type /Pattern /PatternType 2
    /Shading << /ShadingType 2 /ColorSpace /DeviceRGB
      /Coords [0 0 \gradientwidth\space 0]
      /Function << /FunctionType 3 /Domain [0 1]
        /Functions [
          << /FunctionType 2 /Domain [0 1]
            /C0 [\gradientend] /C1 [\gradientend] /N 1 >>
          << /FunctionType 2 /Domain [0 1]
            /C0 [\gradientstart] /C1 [\gradientstart] /N 1 >>
          << /FunctionType 2 /Domain [0 1]
            /C0 [\gradientend] /C1 [\gradientend] /N 1 >>
          << /FunctionType 2 /Domain [0 1]
            /C0 [\gradientstart] /C1 [\gradientstart] /N 1 >>]
        /Bounds [0.25 0.5 0.75] /Encode [0 1 0 1 0 1 0 1] >>
      /Extend [true true] >> >>}%
  \edef\gradientresource{\the\pdflastobj\space 0 R}%
  \setbox\gradientformbox=\hbox{%
    \pdfliteral direct{/Pattern cs /AlignmentGradient scn}%
    \box\gradientbox}%
  \immediate\pdfxform resources{/Pattern << /AlignmentGradient \gradientresource >>}\gradientformbox
  \hbox{\pdfrefxform\pdflastxform}%
\endgroup}
\makeatother

\tikzset{
  header/.style={anchor=north west, text=ink, font=\hd, inner sep=0pt},
  small/.style={text=ink, font=\sm, inner sep=0pt},
  note/.style={rounded corners=0.8mm, line width=0.4pt, draw=border,
               fill=mathbg, inner xsep=3pt, inner ysep=2.2pt, font=\bd,
               text=ink},
  chip/.style={rounded corners=0.7mm, line width=0.35pt, draw=border,
               fill=mathbg, inner xsep=2.6pt, inner ysep=1.8pt, font=\sm,
               text=ink},
  panelbox/.style={rounded corners=1.4mm, line width=0.4pt, draw=border,
                   fill=panelbg},
  flow/.style={-{Stealth[length=1.2mm,width=0.9mm]}, line width=0.4pt,
               draw=rule},
}

% a light tint behind a piece of an equation: \hl{content}
\newcommand{\hl}[1]{\tikz[baseline=(hlx.base)]{%
  \node[inner xsep=1.6pt, inner ysep=1.0pt, rounded corners=0.4mm,
        fill=mathbg] (hlx) {#1};}}

% One rollout, as four tokens.  Each token type is a colour AND a shape, so
% "same string" against "different string" is legible in greyscale and under
% every colour-vision deficiency; minimum sizes are tuned for equal visual
% weight, not equal bounding box (a triangle at equal area looks small).
\tikzset{
  tokglyph/.style={draw=ink!40, line width=0.2pt, inner sep=0pt},
  % shapes, sized for equal visual weight rather than equal bounding box
  % (a triangle at equal area reads small)
  tsquare/.style={tokglyph, rectangle, minimum size=0.072in},
  ttriangle/.style={tokglyph, regular polygon, regular polygon sides=3,
                    minimum size=0.092in},
  tcircle/.style={tokglyph, circle, minimum size=0.080in},
  tdiamond/.style={tokglyph, diamond, aspect=1, minimum size=0.098in},
  % the tokens themselves, named so a rollout reads as what it draws
  bluebox/.style={tsquare, fill=tokblue},
  pinktri/.style={ttriangle, fill=tokpink},
  tealcir/.style={tcircle, fill=tokteal},
  golddia/.style={tdiamond, fill=tokgold},
  greendia/.style={tdiamond, fill=tokgreen},
  yellowtri/.style={ttriangle, fill=tokgold},
  pinkcir/.style={tcircle, fill=tokpink},
}
% \rollout{x-left}{y-centre}{4 token styles}
\newcommand{\rollout}[3]{%
  \foreach \t [count=\i from 0] in {#3} {%
    \node[\t] at ($(#1,#2)+(\i*0.092,0)$) {};}}

\begin{document}
\begin{tikzpicture}[x=1in, y=1in]
% ===========================================================================
% Layout in inches. The top row introduces the metric and an exact-match
% example; the bottom row retains the decomposition and convergence evidence.
% The equation pairs names with symbols; bar ticks repeat only names. Each
% generated part checks that its labels fit; inspect the composite after edits.
% ===========================================================================
\path[use as bounding box] (-0.05,-0.01) rectangle (5.55,3.830);
\path[panelbox] (-0.030,2.825) rectangle (1.530,3.810);
\path[panelbox] ( 1.590,2.825) rectangle (3.620,3.810);
\path[panelbox] ( 3.680,2.825) rectangle (5.530,3.810);
\path[panelbox] (-0.030,0.010) rectangle (2.830,2.760);
\path[panelbox] ( 2.870,1.130) rectangle (5.530,2.760);
\path[panelbox] ( 2.870,0.010) rectangle (5.530,1.070);

% ===================================================================== (a)
% -------------------------- (a.i) four independent calls to the same LLM
\node[header, align=left] at (0.02,3.780)
  {\textbf{(a.i)}~LLM outputs can\\collide or be diverse};
% Each small box contains one independently sampled response pair.
\path[rounded corners=0.7mm, line width=0.35pt, draw=border, fill=mathbg]
  (-0.01,3.175) rectangle (1.51,3.505);
\path[rounded corners=0.7mm, line width=0.35pt, draw=border, fill=mathbg]
  (-0.01,2.840) rectangle (1.51,3.145);
% Every response has its own prompt and LLM box: never branch one call into
% two responses. The repeated LLM denotes the same model in all four rows.
\foreach \ye [count=\i] in {3.420,3.265,3.075,2.920} {
  \node[chip, inner xsep=1.6pt] (prompt-\i) at (0.05,\ye) {$x$};
  \node[chip, inner xsep=1.6pt] (llm-\i) at (0.34,\ye) {LLM};
  \draw[flow] (prompt-\i.east) -- (llm-\i.west);
  \draw[flow] (llm-\i.east) -- (0.665,\ye);}
\rollout{0.74}{3.420}{golddia,tealcir,bluebox,pinktri}
\rollout{0.74}{3.265}{golddia,tealcir,bluebox,pinktri}
\node[text=ink, font=\bd, anchor=center] at (1.105,3.3425) {$=$};
\node[small, anchor=center] at (1.31,3.3425) {collide};

\rollout{0.74}{3.075}{golddia,tealcir,bluebox,pinktri}
\rollout{0.74}{2.920}{bluebox,greendia,pinkcir,yellowtri}
\node[text=ink, font=\bd, anchor=center] at (1.105,2.9975) {$\neq$};
\node[small, anchor=center] at (1.31,2.9975) {diverse};

% ------------------------------------------------- (a.ii) notation
\node[header, align=left, text width=1.90in] at (1.64,3.780)
  {\textbf{(a.ii)}~The notation used in this figure for collision};
\node[note, anchor=north west, font=\eqn] at (1.64,3.470)
  {$C_k(p)\coloneqq\mathbb{E}_{Y,Y'\sim\,p(\cdot\mid x)}\;k(Y,Y')$};
\node[small, anchor=north west, align=flush left, text width=1.90in] at (1.64,3.220)
  {$p$: target; $q$: model; $x$: prompt.\\
   $Y,Y'$: independent responses.\\
   $k$: similarity kernel; $C_k$: mean similarity.};

% ------------------------------------------------ (a) right: what R compares
\node[header, align=left] at (3.73,3.780)
  {\textbf{(a.iii)}~Definition of the\\collision ratio};
\node[anchor=north west, inner sep=0pt, text=ink, font=\eqn] at (3.73,3.440)
  {$R=\dfrac{C_k(q)}{C_k(p)}$};
\node[small, anchor=north west, align=flush left] at (4.57,3.440)
  {$=1$: calibrated\\
   \textcolor{under}{$>1$: under-diverse}\\
   \textcolor{over}{$<1$: over-diverse}};

% ============================================================ row-2 headers
\node[header, align=left, text width=2.76in] at (0.02,2.735)
  {\textbf{(b)}~Model outputs can be over-diverse or under-diverse (mode-collapsed), depending on model and data {\sm(Eq.~(6); Experiment~1)}};
\node[header, align=left, text width=2.56in] at (2.90,2.735)
  {\textbf{(c)}~Fine-tuning calibrates the diversity of the models {\sm(Experiments~2--3)}};

% ================================ the two boxed laws, under their own headers
% The collision identity explains (b). The probability-fit bound has its
% own small panel (d), separate from the empirical training curves in (c).
\node[note, anchor=north west] at (0.02,2.310)
  {$\mathbb{E}[R]=1+\dfrac{1}{C(p_h)}\Big[
    \textcolor{under}{\underbrace{\mathrm{Tr}\,\mathrm{Cov}(q_h)}_{\text{\sm Variance}}}
    +\textcolor{under}{\underbrace{\|b_h\|_2^2}_{\mathclap{\text{\sm Squared bias}}}}
    \;+\;\gradientmath{\underbrace{2p_h^{\!\top}b_h}_{\text{\sm Alignment}}}\Big]$};

% ===================================== (b) LLMs can be over- and under-diverse
\node[anchor=south west, inner sep=0pt] at (0.00,0.19)
  {\includegraphics{panel_sign.pdf}};

% ==================================== (c) fine-tuning calibrates the diversity
\node[anchor=south west, inner sep=0pt] at (2.90,1.11)
  {\includegraphics{panel_close.pdf}};

% ==================================== (d) theoretical probability-fit bound
\node[header, align=left, text width=2.56in] at (2.90,1.040)
  {\textbf{(d)}~A model sufficiently close to optimal under population cross-entropy cannot exhibit arbitrarily miscalibrated diversity {\sm(Theorem~1)}};
\node[note, anchor=north west] at (2.90,0.540)
  {$|C_k(p)-C_k(q)|\le\sqrt{\mathrm{KL}(p\,\|\,q)}$};
\node[small, anchor=north west, align=flush left, text width=2.56in] at (2.92,0.290)
  {$\mathrm{KL}(p\|q)$: Kullback--Leibler divergence,
   equal to excess population cross-entropy.};

\end{tikzpicture}
\end{document}
